% Source: 09_PINN_Layer/holdout_curves.csv.gz, six tall_makeup OP-6 holdouts.
% Pool station-time records, filter on fds_m, then average abs(pinn_m-fds_m).
\reportfigure{\begin{tikzpicture}
\begin{axis}[name=series,width=.50\linewidth,height=4.25cm,xmin=0,xmax=120,
 ymin=.30,ymax=.80,xtick={0,60,120},ytick={.38,.55,.76},ymajorgrids,
 xlabel={时间 (s)},ylabel={净高 (m)},title={远端测点时序},
 legend style={at={(.5,1.01)},anchor=south,legend columns=2,font=\sffamily\scriptsize},
 title style={yshift=12pt,font=\sffamily\small\bfseries}]
\addplot[teal,thin,mark=*,mark size=.65pt] table[col sep=comma,x=time_s,y=fds_m]{data/pinn-curve-example.csv};\addlegendentry{FDS 原始输出}
\addplot[rust,thick,dashed,no marks] table[col sep=comma,x=time_s,y=pinn_m]{data/pinn-curve-example.csv};\addlegendentry{PINN 预测}
\addplot[ink!55,densely dotted,forget plot] coordinates {(0,.38)(120,.38)};
\end{axis}
\begin{axis}[at={(series.east)},anchor=west,xshift=1.90cm,width=.43\linewidth,height=4.25cm,
 xmin=0,xmax=13,xtick={0,5,10},ymin=.48,ymax=3.52,ytick={1,2,3},
 yticklabels={$h\leq0.71$ m,$h<0.76$ m,全程},xmajorgrids,
 xlabel={净高 MAE (cm)},title={六个 OP-6 留出算例},
 title style={yshift=12pt,font=\sffamily\small\bfseries}]
\addplot[xbar,bar width=5pt,draw=none,fill=teal,
 nodes near coords={\pgfmathprintnumber[fixed,precision=2]{\pgfplotspointmeta}},
 point meta=x,nodes near coords align={horizontal},
 every node near coord/.append style={font=\sffamily\scriptsize,text=ink}]
 coordinates {(1.68,3)(4.75,2)(10.09,1)};
\end{axis}
\end{tikzpicture}}{图 6.\quad 左图为曲线 MAE 最接近组均值的算例（1.28~m$^3$/s，烟炱产率 0.02，垂壁深度 0.24~m）在 $x=6.1$~m 的时序，水平点线为 0.38~m 判据。右图按 FDS 净高 $h$ 筛选六例的测点记录，由上至下分别含 11\,568、2\,001 和 553 条；第三组为第二组的子集。}
